\documentclass[a4paper,12pt]{article}
\usepackage[T2A]{fontenc}
\usepackage[utf8]{inputenc}
\usepackage[russian]{babel}
\usepackage{amsmath,amssymb,mathtools}
\usepackage{helvet}
\renewcommand{\familydefault}{\sfdefault}
\usepackage{geometry}
\geometry{left=22mm,right=22mm,top=23mm,bottom=24mm}
\usepackage{microtype,icomma}
\usepackage{xcolor}
\usepackage{tikz}
\usetikzlibrary{arrows.meta,positioning,shapes.geometric,calc}
\usepackage{booktabs,tabularx,array}
\usepackage{enumitem,parskip,fancyhdr,hyperref}
\usepackage{setspace}
\definecolor{accent}{HTML}{008080}
\definecolor{darkaccent}{HTML}{005354}
\definecolor{textgray}{HTML}{343434}
\definecolor{muted}{HTML}{6B7474}
\definecolor{rulecolor}{HTML}{CADCDC}
\hypersetup{colorlinks=true,urlcolor=darkaccent,linkcolor=darkaccent}
\setstretch{1.10}
\setlength{\parindent}{0pt}
\setlength{\parskip}{4pt}
\setcounter{secnumdepth}{0}
\setlist[itemize]{leftmargin=6mm,itemsep=3pt,topsep=3pt}
\setlist[enumerate]{leftmargin=8mm,itemsep=6pt,topsep=4pt}
\renewcommand{\arraystretch}{1.22}
\setlength{\headheight}{25pt}
\pagestyle{fancy}
\fancyhf{}
\fancyhead[L]{\small\color{darkaccent}Задачи на движение}
\fancyhead[R]{\small\color{muted}Ярослав Гаврилов}
\fancyfoot[C]{\small\color{muted}\thepage}
\renewcommand{\headrulewidth}{0.35pt}
\renewcommand{\headrule}{\hbox to\headwidth{\color{rulecolor}\leaders\hrule height \headrulewidth\hfill}}
\renewcommand{\footrulewidth}{0pt}
\newcommand{\sectiontitle}[1]{\section*{\color{darkaccent}#1}\vspace{-2mm}}
\newcommand{\key}[1]{\textcolor{darkaccent}{\textbf{#1}}}
\newcommand{\smallnote}[1]{{\small\color{muted}#1}}
\newcommand{\exlabel}[1]{\textbf{\color{darkaccent}Пример.} #1}
\newcommand{\vrel}{v_{\text{отн}}}
\newcommand{\vavg}{v_{\text{ср}}}
\newcommand{\stot}{S_{\text{общ}}}
\newcommand{\ttot}{t_{\text{общ}}}
\begin{document}
\begin{titlepage}
\thispagestyle{empty}
\begin{center}
\vspace*{15mm}
{\Large\color{darkaccent}\textsc{Математика}\par}
\vspace{17mm}
{\fontsize{25}{30}\selectfont Чек-лист\par}
\vspace{11mm}
{\fontsize{20}{25}\selectfont\color{accent}Задачи на движение:\par}
\vspace{2mm}
{\fontsize{19}{24}\selectfont\color{accent}средняя скорость и\par}
{\fontsize{19}{24}\selectfont\color{accent}протяжённые тела\par}
\vspace{15mm}
{\large Подготовка к ЕГЭ по профильной математике\par}
\vfill
{\large Ярослав Гаврилов\par}
\vspace{12mm}
{\large 23 ноября 2026 г.\par}
\vspace{11mm}
{\normalsize Лёвин Артём Александрович\par}
\vspace{2mm}
{\normalsize эксклюзивно для Ярослава Гаврилова\par}
\vspace{10mm}
\begin{tikzpicture}[x=1cm,y=1cm]
\draw[accent,line width=0.85pt] (0,0) .. controls (1.4,0.65) and (2.8,-0.5) .. (4.25,0)
 .. controls (5.65,0.48) and (7.4,-0.5) .. (9,0.02);
\fill[accent] (4.25,0) circle (1.1pt);
\end{tikzpicture}
\end{center}
\end{titlepage}
\setcounter{page}{2}

\sectiontitle{1. Общая основа: скорость, путь, время}
Для любого участка равномерного движения используйте три взаимосвязанные величины: скорость $v$, время $t$ и пройденный путь $S$.
\[
\boxed{S=vt,\qquad t=\frac{S}{v},\qquad v=\frac{S}{t}.}
\]
Удобно заполнять таблицу: по её строкам видно, какой путь пройден с какой скоростью и сколько времени на это потребовалось.
\begin{center}
\begin{tabular}{@{}p{39mm}p{37mm}p{38mm}p{33mm}@{}}
\toprule
Участок & Скорость & Время & Путь \\
\midrule
Первый & $v_1$ & $t_1=S_1/v_1$ & $S_1$ \\
Второй & $v_2$ & $t_2=S_2/v_2$ & $S_2$ \\
\bottomrule
\end{tabular}
\end{center}
\key{Первый вопрос к условию:} делятся ли пополам время, расстояние или ни то ни другое? От этого зависит правильная запись таблицы.

\sectiontitle{2. Движение судна по реке}
Обозначим $x$ собственную скорость судна, то есть скорость в стоячей воде, а $c$ --- скорость течения.
\[
\boxed{v_{\text{по течению}}=x+c,\qquad
v_{\text{против течения}}=x-c,\qquad x>c.}
\]
Если расстояние в каждом направлении равно $S$, времена равны $S/(x+c)$ и $S/(x-c)$. Когда движение против течения заняло на $\Delta t$ часов больше, составляем уравнение
\[
\boxed{\frac{S}{x-c}-\frac{S}{x+c}=\Delta t.}
\]
\exlabel{Судно проходит $24$ км против течения на $2$ часа дольше, чем $24$ км по течению. Скорость течения $3$ км/ч. Найдём собственную скорость судна.}
\[
\frac{24}{x-3}-\frac{24}{x+3}=2,\qquad x>3.
\]
Приводим к общему знаменателю, \emph{не теряя скобок}:
\[
\frac{24(x+3)-24(x-3)}{(x-3)(x+3)}=2
\quad\Longrightarrow\quad
\frac{144}{x^2-9}=2.
\]
\[
x^2-9=72,\qquad x^2=81,\qquad x=\pm9.
\]
Отрицательная скорость невозможна, а также должно выполняться $x>3$. Поэтому \key{ответ: $9$ км/ч}. Проверка: $24/6=4$ ч против течения и $24/12=2$ ч по течению; разница действительно $2$ ч.

\clearpage
\sectiontitle{3. Средняя скорость: главный принцип}
Средняя скорость по всему маршруту --- не произвольное среднее арифметическое скоростей. Она равна отношению \emph{всего пути} ко \emph{всему времени}:
\[
\boxed{\vavg=\frac{S_1+S_2+\dots+S_n}{t_1+t_2+\dots+t_n}.}
\]
\key{Порядок действий:} разделите движение на участки; для каждого определите путь и время; сложите пути, сложите времена; разделите одну сумму на другую.

\sectiontitle{4. Две половины \emph{времени}}
\exlabel{Автомобиль первую половину общего времени ехал со скоростью $84$ км/ч, вторую --- $56$ км/ч.}
Обозначим продолжительность всей поездки через $t$ часов. Тогда каждый участок занял $t/2$ часов, а пройденные пути различаются.
\begin{center}
\begin{tabular}{@{}p{42mm}p{36mm}p{40mm}p{33mm}@{}}
\toprule
Часть поездки & Скорость & Время & Путь \\
\midrule
Первая половина времени & $84$ & $t/2$ & $42t$ \\
Вторая половина времени & $56$ & $t/2$ & $28t$ \\
\bottomrule
\end{tabular}
\end{center}
\[
\vavg=\frac{42t+28t}{t/2+t/2}=\frac{70t}{t}=
\boxed{70\ \text{км/ч}}.
\]
Если \emph{времена равны}, можно пользоваться краткой формулой
\[
\boxed{\vavg=\frac{v_1+v_2}{2}.}
\]

\sectiontitle{5. Две половины \emph{пути}}
\exlabel{Автомобиль первую половину маршрута прошёл со скоростью $56$ км/ч, вторую --- $84$ км/ч.}
Весь путь равен $S$ км. Каждый участок имеет длину $S/2$, однако время \emph{разное}: медленная половина длится дольше. Поэтому результат будет меньше $70$ км/ч из предыдущего примера.
\begin{center}
\begin{tabular}{@{}p{42mm}p{36mm}p{40mm}p{33mm}@{}}
\toprule
Часть поездки & Скорость & Время & Путь \\
\midrule
Первая половина пути & $56$ & $S/112$ & $S/2$ \\
Вторая половина пути & $84$ & $S/168$ & $S/2$ \\
\bottomrule
\end{tabular}
\end{center}
\[
\vavg=\frac{S/2+S/2}{S/112+S/168}
=\frac{1}{\frac1{112}+\frac1{168}}
=\frac{1}{\frac{3+2}{336}}
=\boxed{67{,}2\ \text{км/ч}}.
\]
\vspace{-2mm}\noindent\smallnote{При равных расстояниях можно использовать $\vavg=2v_1v_2/(v_1+v_2)$.}
\clearpage
\sectiontitle{6. Протяжённые тела: поезд, мост, другой поезд}
Когда поезд проходит мимо столба, пешехода, моста или другого поезда, учитывайте \emph{длину движущегося тела}. Время полного прохождения рассчитывается по относительному движению:
\[
\boxed{t=\frac{S_{\text{отн}}}{\vrel}},\qquad
S_{\text{отн}}=L_1+L_2
\quad\text{(если начальных и конечных промежутков нет).}
\]
Здесь $L_1$ и $L_2$ --- длины объектов вдоль направления движения; длину столба и пешехода в типовых школьных задачах принимают равной нулю относительно длины поезда.

\begin{center}
\begin{tabularx}{\linewidth}{@{}p{50mm}X@{}}
\toprule
Как движутся объекты & Относительная скорость \\
\midrule
Навстречу друг другу & $\vrel=v_1+v_2$ \\
В одном направлении; один догоняет другой & $\vrel=v_{\text{быстрого}}-v_{\text{медленного}}$ \\
Один объект неподвижен & $\vrel=v_{\text{движущегося}}$ \\
\bottomrule
\end{tabularx}
\end{center}
\key{Уточнение смысла времени.} В задаче «поезд проходит мимо человека» путь относительно человека равен длине поезда. В задаче «поезд полностью проходит мимо другого поезда» учитывается сумма длин \emph{обоих} поездов. Эти два условия нельзя подменять друг другом.

\sectiontitle{7. Согласование единиц измерения}
Если скорость записана в км/ч, удобно выразить длину в километрах, а время --- в часах. Если скорость в м/с, используйте метры и секунды.
\[
1\ \text{км}=1000\ \text{м},\quad
1\ \text{ч}=60\ \text{мин}=3600\ \text{с},\quad
1\ \text{м/с}=3{,}6\ \text{км/ч}.
\]
\[
30\ \text{с}=\frac{30}{3600}\ \text{ч},\qquad
18\ \text{мин}=\frac{18}{60}\ \text{ч},\qquad
600\ \text{м}=0{,}6\ \text{км}.
\]
\exlabel{Поезд со скоростью $60$ км/ч проходит мимо неподвижного столба за $30$ секунд. Найдём длину поезда.}
\[
L=v\,t=60\cdot\frac{30}{3600}=0{,}5\ \text{км}
=\boxed{500\ \text{м}}.
\]
\exlabel{Поезд движется со скоростью $76$ км/ч навстречу пешеходу, который идёт со скоростью $4$ км/ч. Прохождение поезда мимо человека длится $18$ секунд.}
\[
\vrel=76+4=80\ \text{км/ч},\qquad
L=80\cdot\frac{18}{3600}=0{,}4\ \text{км}
=\boxed{400\ \text{м}}.
\]

\clearpage
\sectiontitle{8. Обгон с начальным и конечным промежутками}
Если быстрый объект сначала находится позади медленного, а затем полностью обгоняет его и уходит вперёд, одного сложения длин недостаточно. Быстрому объекту нужно в относительном движении преодолеть:
\[
\boxed{S_{\text{отн}}=d_{\text{до}}+L_1+L_2+d_{\text{после}}.}
\]
Здесь $d_{\text{до}}$ --- расстояние между ближайшими концами объектов до обгона, $d_{\text{после}}$ --- расстояние между ближайшими концами после обгона. При движении в одном направлении:
\[
\boxed{t=\frac{d_{\text{до}}+L_1+L_2+d_{\text{после}}}
{v_{\text{быстрого}}-v_{\text{медленного}}}.}
\]
\exlabel{Две баржи движутся в одном направлении. Их длины $60$ и $40$ метров. Сначала между кормой передней баржи и носом задней $200$ метров. Через $18$ минут задняя баржа полностью обогнала переднюю, и расстояние от её кормы до носа обогнанной баржи стало $300$ метров. На сколько км/ч скорость задней баржи больше?}
\[
S_{\text{отн}}=60+40+200+300=600\ \text{м}=0{,}6\ \text{км},
\qquad t=\frac{18}{60}=0{,}3\ \text{ч}.
\]
Нужна именно \emph{разность скоростей} (обозначим её $\Delta v$), поэтому отдельно скорости каждой баржи искать не требуется:
\[
\Delta v=\frac{0{,}6}{0{,}3}=\boxed{2\ \text{км/ч}}.
\]
\key{Правильный геометрический смысл:} сумма двух длин и двух промежутков --- это полный путь, на который задняя баржа должна «выиграть» у передней.

\sectiontitle{9. Ошибки, которых важно избегать}
\begin{itemize}
\item \textbf{Путать половину времени с половиной пути.} Равные времена не означают равных расстояний и наоборот.
\item \textbf{Усреднять скорости без проверки условий.} Формула $(v_1+v_2)/2$ подходит при равных временах, но обычно не при равных расстояниях.
\item \textbf{Неверно выбирать знак относительной скорости.} При сближении навстречу скорости складываются; при догонянии --- вычитаются.
\item \textbf{Смешивать километры, метры, минуты и секунды.} Перед подстановкой согласуйте все единицы.
\item \textbf{Забывать вторую длину или промежуток.} Сначала определите, какой именно объект проходит мимо какого и что означает начало и конец измеряемого времени.
\item \textbf{Сокращать алгебраические дроби без ОДЗ.} В задаче о течении обязательно $x>c$; при раскрытии скобок в числителе учитывайте знак минус перед второй дробью.
\end{itemize}

\clearpage
\thispagestyle{fancy}
\begin{center}
{\Large\bfseries\color{darkaccent}Алгоритм: прохождение протяжённых тел}\par
\smallnote{Определяем относительный путь и относительную скорость.}
\end{center}
\vspace{8mm}
\begin{center}
\begin{tikzpicture}[
  node distance=8mm,
  startstop/.style={ellipse,draw=accent,line width=.8pt,align=center,minimum width=50mm,minimum height=9mm,font=\small},
  process/.style={rectangle,rounded corners=2mm,draw=rulecolor,line width=.85pt,align=center,text width=112mm,minimum height=13mm,inner sep=3mm,font=\small},
  choice/.style={diamond,draw=accent,aspect=2.4,align=center,text width=35mm,inner sep=1.4mm,font=\small},
  branch/.style={rectangle,rounded corners=2mm,draw=rulecolor,align=center,text width=49mm,minimum height=13mm,inner sep=2mm,font=\small},
  arr/.style={-{Latex[length=2.6mm]},line width=.8pt,color=accent}
]
\node[startstop] (start) at (0,0) {Начало};
\node[process] (read) at (0,-1.85) {Определите, что проходит мимо чего и какой момент считается началом и концом движения};
\node[process] (path) at (0,-4.10) {Найдите относительный путь: сумма нужных длин и, если они есть, начального и конечного промежутков};
\node[choice] (test) at (0,-6.65) {Объекты движутся навстречу?};
\node[branch] (opposite) at (-3.65,-9.05) {$\vrel=v_1+v_2$};
\node[branch] (same) at (3.65,-9.05) {Вдогонку или мимо столба \\ $\vrel=v_{\text{быстр.}}-v_{\text{медл.}}$};
\node[process] (units) at (0,-11.23) {Согласуйте единицы: км и часы либо метры и секунды};
\node[process] (calc) at (0,-13.45) {Используйте $t=S_{\text{отн}}/\vrel$ и выразите искомое время, длину или разность скоростей};
\node[startstop] (finish) at (0,-15.4) {Проверьте ответ и размерность};
\draw[arr] (start) -- (read);
\draw[arr] (read) -- (path);
\draw[arr] (path) -- (test);
\draw[arr] (test.west) -| node[pos=.23,above,font=\small]{да} (opposite.north);
\draw[arr] (test.east) -| node[pos=.23,above,font=\small]{нет} (same.north);
\draw[arr] (opposite.south) -- ++(0,-.57) -| ([xshift=-15mm]units.north);
\draw[arr] (same.south) -- ++(0,-.57) -| ([xshift=15mm]units.north);
\draw[arr] (units) -- (calc);
\draw[arr] (calc) -- (finish);
\end{tikzpicture}
\end{center}
\clearpage

\sectiontitle{Тренировочный блок}
\smallnote{Путь А: 10 задач. Выполняйте последовательно, от базовых ситуаций к составным. Все ответы указывайте с единицами измерения; при необходимости сначала составляйте таблицу «скорость --- время --- путь».}
\vspace{3mm}
\begin{enumerate}[label=\textcolor{darkaccent}{\textbf{\arabic*.}},leftmargin=9mm,itemsep=11pt]
\item \textbf{Две половины времени.} Автомобиль выехал из города на загородную дорогу и двигался без остановок. Первую половину общего времени он поддерживал скорость $60$ км/ч, а во вторую половину, когда дорога стала свободнее, --- $90$ км/ч. Найдите среднюю скорость автомобиля за всю поездку в км/ч.

\item \textbf{Две половины маршрута.} Поездка на автомобиле составила $180$ км. Первые $90$ км водитель проехал по обычной дороге со скоростью $60$ км/ч, оставшиеся $90$ км --- по скоростной трассе со скоростью $90$ км/ч. Остановок не было. Определите среднюю скорость на всём маршруте в км/ч и сравните её с результатом предыдущей задачи.

\item \textbf{Три участка с разной продолжительностью.} Междугородний автобус ехал четыре часа без остановок. В первый час его скорость была $48$ км/ч, в следующие два часа --- $72$ км/ч, а в последний час --- $96$ км/ч. Какова средняя скорость автобуса за все четыре часа? Ответ дайте в км/ч.

\item \textbf{Поезд и придорожный столб.} Пассажирский поезд движется равномерно со скоростью $72$ км/ч. От момента, когда его локомотив поравнялся с неподвижным сигнальным столбом, до момента прохождения мимо столба последнего вагона прошло $25$ секунд. Найдите полную длину поезда в метрах.

\item \textbf{Поезд и пешеход навстречу.} Вдоль железнодорожного полотна идёт пешеход со скоростью $5{,}4$ км/ч. Навстречу ему по соседнему пути движется поезд со скоростью $72$ км/ч. С того момента, когда локомотив поравнялся с пешеходом, до прохождения последнего вагона прошло $20$ секунд. Определите длину поезда в метрах, считая длину пешехода пренебрежимо малой.

\pagebreak
\item \textbf{Поезд обгоняет пешехода.} Поезд со скоростью $72$ км/ч обгоняет пешехода, который идёт параллельно путям в том же направлении со скоростью $3{,}6$ км/ч. Полное прохождение поезда мимо пешехода занимает $30$ секунд. Найдите длину состава в метрах. Объяснение выбора разности скоростей приведите в решении.

\item \textbf{Два встречных поезда.} По параллельным путям навстречу друг другу движутся пассажирский поезд длиной $240$ м со скоростью $54$ км/ч и грузовой поезд длиной $360$ м со скоростью $90$ км/ч. Сколько секунд пройдёт от встречи локомотивов до момента, когда последние вагоны полностью разойдутся?

\item \textbf{Обгон одного поезда другим.} На соседних путях в одном направлении движутся два поезда: быстрый длиной $180$ м со скоростью $108$ км/ч и медленный длиной $320$ м со скоростью $72$ км/ч. Сколько секунд длится полный обгон, если отсчёт начинается в момент, когда локомотив быстрого поезда поравнялся с последним вагоном медленного, и заканчивается, когда последний вагон быстрого поезда обогнал локомотив медленного?

\item \textbf{Катер на реке.} Катер прошёл $24$ км против течения реки, а затем вернулся тем же путём. Скорость течения равна $3$ км/ч. На участок против течения катер затратил на $2$ часа больше, чем на обратный участок. Найдите собственную скорость катера в км/ч; учтите, что она должна быть больше скорости течения.

\item \textbf{Две баржи и полный обгон.} По прямому каналу в одном направлении движутся две баржи. Длина передней баржи $70$ м, длина более быстрой задней --- $50$ м. В начале наблюдения нос задней баржи находился в $180$ м позади кормы передней. Через $12$ минут задняя баржа полностью обогнала переднюю, и расстояние от кормы задней до носа передней стало $300$ м. На сколько километров в час скорость задней баржи больше скорости передней?
\end{enumerate}

\clearpage
\sectiontitle{Ответы к тренировочному блоку}
\smallnote{Проверьте результат после самостоятельного решения. Ответы не заменяют запись формулы, выбора относительной скорости и согласования единиц.}
\vspace{6mm}
\begin{center}
\begin{tabularx}{.84\linewidth}{@{}>{\centering\arraybackslash}p{16mm}X@{}}
\toprule
\textbf{№} & \textbf{Ответ} \\
\midrule
1 & $75$ км/ч \\
2 & $72$ км/ч \\
3 & $72$ км/ч \\
4 & $500$ м \\
5 & $430$ м \\
6 & $570$ м \\
7 & $15$ с \\
8 & $50$ с \\
9 & $9$ км/ч \\
10 & $3$ км/ч \\
\bottomrule
\end{tabularx}
\end{center}
\vspace{12mm}
\sectiontitle{Итоговая памятка}
\begin{tabularx}{\linewidth}{@{}p{62mm}X@{}}
\toprule
\textbf{Сигнал в условии} & \textbf{Первое действие} \\
\midrule
Половина времени & Запишите одинаковые времена; путь $S_i=v_it_i$. \\
Половина пути & Запишите одинаковые пути; время $t_i=S_i/v_i$. \\
Средняя скорость & Сложите все пути и все времена отдельно. \\
Поезд мимо столба или человека & Определите длину поезда и относительную скорость. \\
Встречное движение & Сложите скорости объектов. \\
Догоняющее движение & Вычтите меньшую скорость из большей. \\
Полный обгон с промежутками & К длинам прибавьте промежутки до и после. \\
Движение по реке & Используйте $x-c$ и $x+c$, учитывайте $x>c$. \\
\bottomrule
\end{tabularx}
\vfill
\begin{center}
\smallnote{Лёвин Артём Александрович --- эксклюзивно для Ярослава Гаврилова}
\end{center}
\end{document}
